\subsection{Lagwise packet--HLZ bridge for the log-only source}
\label{sec:lagwise-packet-hlz}

The Fourier-diagonalization sentence following
\eqref{eq:pre-HLZ-lag-factorization} is not used in the revised principal
route.  In particular, we do not identify the Fourier multiplier of a lag
convolution with a one-sided packet-prefix multiplier.  Instead we retain the
actual continuous lag $d=y-x$ and apply the scalar HLZ diagonalization
uniformly at fixed $d$.

\begin{proposition}[Lagwise packet--HLZ bridge]
\label{prop:lagwise-packet-HLZ}
Restrict to the invariant source $h=k=1$, one retained height block, and the
stable flat packet core.  Put
\[
X(\tau)=\frac{\tau}{2\pi},\qquad
\lambda(\tau)=\frac{\log X(\tau)}L,
\]
and, for $|d|\le W_T$, define
\begin{equation}
\boxed{
\rho_\tau(d):=\lambda(\tau)-\frac{|d|}{L}.
}
\label{eq:lagwise-rho}
\end{equation}
For the logarithmic source $H_{\log}$ of \eqref{eq:Hlog}, the principal
stationary packet form after the finite three-$Z$ Cauchy residues is
\begin{align}
\mathcal Q_{\log,T}(v,w;p)
={}&\iint
 (\mathcal A_Tv)(x)\overline{(\mathcal A_Tw)(y)}
 \mathcal G^{\log}_{\lambda(\tau),\rho_\tau(y-x)}(p)\,dx\,dy
\notag\\
&+\mathcal R^{\rm HLZ}_{\log,T}(v,w;p),
\label{eq:lagwise-HLZ-form}
\end{align}
where
\begin{equation}
\boxed{
\mathcal G^{\log}_{\lambda,\rho}(p)
=e^{\lambda p}\left[
-2E_\rho(p)^2
+e^{2\lambda}E_\rho(p+2)E_\rho(p-2)
\right],
\qquad
E_\rho(q)=\int_0^\rho e^{-qt}\,dt.
}
\label{eq:logonly-local-profile}
\end{equation}
For every prescribed $A>0$, after the fixed Cauchy and curvature residues have
been taken, the scalar error profile and its first two $\rho$-derivatives are
$O_A(L^{-A})$ up to a fixed polylogarithmic factor, uniformly in
$|d|\le W_T$ and in the retained height blocks.
\end{proposition}

\begin{proof}
For basis packets $e_r,e_s$, the exact polarization gives
\begin{equation}
\mathcal B_{e_r,e_s}(c+it)
=\iint\phi_L(x)\phi_L(y)e^{-i\tau_sx+i\tau_ry}
 e^{\sigma_0(x-y)}e^{it(x-y)}\,dx\,dy,
\label{eq:lagwise-polarization-recall}
\end{equation}
with $\sigma_0=c-1/2$.  For one Dirichlet monomial $\xi^{-c-it}$, put
$d=y-x$.  Then exactly
\[
e^{\sigma_0(x-y)}e^{it(x-y)}\xi^{-c-it}
=e^{d/2}(\xi e^d)^{-c-it}.
\]
Hence the stationary dilation parameter of Section~\ref{sec:stationary} is
$w=d=y-x$; it is not a fixed number attached to either packet label.
The principal stationary calculation depends on the packet variables only
through this lag, exactly as in the first part of
Lemma~\ref{lem:pre-HLZ-packet-separation}.

By Lemma~\ref{lem:HLZ-local-height-rescaling}, the diagonal Mellin factor may be
written at local scale as
\[
\left(\frac{X(\tau)^2}{\ell^2}\right)^u
\]
without changing the $u=0$ residue and with the same arbitrary-log contour
error.  Proposition~\ref{prop:X-tagging} gives the two oriented stationary
cutoffs
\[
\ell\le X(\tau)e^{-d},\qquad
\ell\le X(\tau)e^d,
\]
so their common range is
\begin{equation}
\boxed{
\ell\le X(\tau)e^{-|d|}.
}
\label{eq:lagwise-common-cutoff}
\end{equation}
Writing $t=\log\ell/L$ gives precisely
$0\le t\le\rho_\tau(d)$.  On the stable core
$|d|\le W_T=L-A_0\log L$, so $\rho_\tau(d)$ lies in a fixed compact interval
and is positive for all sufficiently large $T$.

Fix $d$, $p$ and the separated Cauchy variable $\mathbf z$.  After the leading
stationary factors have been removed from the two packet legs, the remaining
scalar Mellin weight has the form
\begin{equation}
W_d(u,\mathbf z,p)
=\int_0^{\rho_\tau(d)}
 a_{d,\mathbf z,p,T}(t)e^{Lut}\,dt,
\label{eq:lagwise-Mellin-weight}
\end{equation}
or a fixed finite product of such factors.  Lemma~\ref{lem:stationary-symbol}
controls the fixed stationary and curvature derivatives of the amplitude by a
fixed power of $L$.  The separated shifts $|z_j|\asymp L^{-1}$ contribute only
fixed polylogarithmic losses.  Moving-endpoint derivatives of
\eqref{eq:lagwise-Mellin-weight} produce only the endpoint jets covered by
Lemma~\ref{lem:smooth-mellin-class}.  Thus $W_d$ belongs to that Mellin class
uniformly for $|d|\le W_T$.

Lemma~\ref{lem:packet-HLZ-diagonal} can therefore be applied pointwise in $d$.
Its arbitrary-log error is uniform after the first two endpoint and curvature
derivatives.  The finite three-$Z$ residues then give the two logarithmic
faces.  The coefficient $-2P$ contributes
\[
-2e^{\lambda p}E_\rho(p)^2,
\]
while $P(s-2/L)$ changes the two diagonal fibre exponents to $p+2$ and $p-2$
and contributes the explicit archimedean factor
$X^{2/L}=e^{2\lambda}$, giving
\[
e^{\lambda p}e^{2\lambda}E_\rho(p+2)E_\rho(p-2).
\]
Their sum is \eqref{eq:logonly-local-profile}.

Finally, on the original absolutely convergent scalar contour the $x,y$
variables range over the compact packet support and the scalar $u$-integral has
Gaussian vertical decay.  Fubini is therefore legitimate before the scalar
contour is moved.  After the move, the uniform error bound just obtained may be
integrated against the same two packet legs.  This proves
\eqref{eq:lagwise-HLZ-form} and the stated differentiated error estimate.
\end{proof}

\begin{corollary}[Frozen log-only lag kernel]
\label{cor:frozen-logonly-lag-kernel}
Let
\begin{equation}
\boxed{
G_{\log}(\rho):=[p^2]\,
 e^p\left[
 -2E_\rho(p)^2+e^2E_\rho(p+2)E_\rho(p-2)
 \right].
}
\label{eq:Glog-from-source}
\end{equation}
After freezing $\lambda(\tau)=1$ and rescaling $r=x/L$, $s=y/L$, the principal
operator is the pullback of the genuine lag kernel
\begin{equation}
\boxed{
K_{\log}(r,s)=G_{\log}(1-|r-s|).
}
\label{eq:Klog-lag-kernel}
\end{equation}
The difference between the finite-height profile
$G_{\log,\lambda(\tau)}(\lambda(\tau)-|x-y|/L)$ and
\eqref{eq:Klog-lag-kernel} has its first two normalized lag derivatives
$O(L^{-1})$ and belongs to the finite-parameter regular class.
\end{corollary}

\begin{proof}
The map $(\lambda,\rho,p)\mapsto\mathcal G^{\log}_{\lambda,\rho}(p)$ is analytic
on a fixed neighbourhood of the retained parameter set and
$\lambda(\tau)=1+O(L^{-1})$.  The mean-value theorem gives the claimed
$O(L^{-1})$ estimates through two $\rho$-derivatives.  Substitution of
$d=y-x$ in Proposition~\ref{prop:lagwise-packet-HLZ} gives the lag kernel.
\end{proof}

\begin{remark}[Replacement of the one-sided bridge]
\label{rem:one-sided-bridge-superseded}
For the proof of Theorem~\ref{thm:main}, Proposition~\ref{prop:lagwise-packet-HLZ}
replaces the identification step in Lemmas~\ref{lem:aux-shift-rigidity-packet-map},
\ref{lem:pre-HLZ-packet-separation}, and
\ref{lem:model-HLZ-Fubini-lift} that attempted to pass from the lag convolution
to a one-sided common-prefix multiplier.  The exact lag factorization
\eqref{eq:pre-HLZ-lag-factorization} itself remains valid; only the subsequent
operator identification is discarded.  In particular the logarithmic
principal source is not sent through the compressed-shift argument of
Section~\ref{sec:edge-rank}.
\end{remark}